% Part: methods
% Chapter: proofs
% Section: using-definitions

\documentclass[../../../include/open-logic-section]{subfiles}

\begin{document}

\olfileid[tr]{mod}{prf}{def}

\olsection{Tanımları Kullanmak}

Kanıtta kullanılabilecek bütün tanımları bilmeniz ve onları doğru
uygulayabilmeniz gerektiğini belirtmiştik. Bu gerçekten önemli bir
noktadır ve biraz daha ayrıntılı incelenmeye değer. Tanımlar, özellik
ve bağıntıları daha kısa biçimde anlatabilmemiz için onları
kısaltmakta kullanılır. Ortaya konan kısaltmaya \emph{tanımlanan}, onun
kısalttığı şeye ise \emph{tanımlayan} denir. Kanıtlarda çoğu kez
tanımlananın nasıl ortaya konduğuna geri dönmek zorunda kalırız; çünkü
kanıtımızı ilerletmek için tanımlayanın (tanımlanan terimin kısalttığı
uzun biçimin) mantıksal yapısından yararlanmamız gerekir. Tanımları
açarak mantıksal işleyişin özüne indiğinizden emin olursunuz.

Bir örnekle başlayalım. Aşağıdakini kanıtlamak istediğinizi düşünün:

\begin{prop}
Herhangi $A$ ve $B$ kümeleri için $A \cup B = B \cup A$.
\end{prop}

Kanıta başlayabilmek için bile iki kümenin özdeş olmasının ne anlama
geldiğini, yani bu eşitlikteki ``$=$'' işaretinin kümeler için ne
anlama geldiğini bilmemiz gerekir. Kümeler, !!{element}s bakımından aynı
olduklarında özdeş olarak tanımlanır. Dolayısıyla açmamız gereken tanım
şudur:

\begin{defn}
$A$ ve $B$ kümeleri
\emph{özdeştir}, $A = B$, ancak ve ancak $A$'daki her !!{element},
$B$'de de bir !!a{element} olarak bulunduğunda ve bunun tersi geçerli
olduğunda.
\end{defn}

Bu tanım, $A$ ve~$B$'yi keyfî kümelerin yer tutucuları olarak kullanır.
Tanımladığı şey---\emph{tanımlanan}---``$A = B$'' ifadesidir ve bunu,
$A = B$'nin doğru olduğu koşulu vererek yapar. Bu koşul---``$A$'daki her
!!{element}, $B$'de de bir !!a{element} olarak bulunur ve bunun tersi de
geçerlidir''---\emph{tanımlayan}dır.
\footnote{Bu özel durumda---ve oldukça kafa karıştırıcı biçimde!---
  $A = B$ olduğunda $A$ ve $B$ kümeleri, sol ve sağ tarafta farklı
  harfler kullansak da, tek ve aynı kümedir. Ancak bu kümenin
  belirlenme yolları farklı olabilir ve tanımı önemsiz olmaktan
  çıkaran da budur.} Tanım, $A = B$'nin ancak ve ancak bu koşul
sağlandığında doğru olduğunu belirtir.

Tanımı uygularken, tanımdaki $A$ ve $B$'yi ele aldığınız durumla
eşleştirmeniz gerekir. Bizim durumumuzda bu, $A \cup B = B \cup A$'nın
doğru olması için her $z \in A \cup B$'nin aynı zamanda $B \cup A$'da
olması ve bunun tersinin de geçerli olması gerektiği anlamına gelir.
Önermedeki $A \cup B$ ifadesi tanımdaki~$A$'nın, $B \cup A$ ise
tanımdaki~$B$'nin rolünü oynar. $A$ ve $B$ hem tanımda hem de
kanıtladığımız önermenin ifadesinde, fakat farklı amaçlarla
kullanıldığından, ikisini karıştırmadığınızdan emin olmalısınız.
Örneğin önermeyi, $A$'daki her !!{element} için $B$'de de bir
!!a{element} bulunduğunu ve bunun tersini göstererek kanıtlayabileceğinizi
düşünmek hata olurdu---bu, $A = B$'yi gösterirdi; $A \cup B = B \cup A$'yı
değil. (Üstelik $A$ ve $B$ herhangi iki küme olabileceğinden pek
ilerleyemezdiniz; çünkü $A$ ve~$B$ hakkında hiçbir şey varsayılmamışsa
pekâlâ farklı kümeler olabilirler.)

Kanıtta birleşim gibi küme kuramsal kavramlarla uğraşıyoruz; bu nedenle
kanıtın nasıl ilerlemesi gerektiğini anlamak için $\cup$ simgesinin
anlamını da bilmemiz gerekir. Bazen bir tanımı açmak, açılması gereken
başka tanımlar doğurur. Örneğin $A \cup B$, $\Setabs{z}{z \in A
  \text{ ya da } z \in B}$ olarak tanımlanır. Dolayısıyla $x \in A \cup
B$ olduğunu kanıtlamak istiyorsanız, $\cup$ tanımını açınca $x \in
\Setabs{z}{z \in A \text{ ya da } z \in B}$ olduğunu kanıtlamanız
gerektiğini görürsünüz. Şimdi ayrıca $x \in
\Setabs{z}{\dots z\dots}$ olmasının ancak ve ancak $\dots x\dots$
olduğunda geçerli olduğunu hatırlamanız gerekir. Öyleyse
$\Setabs{z}{\dots z \dots}$ gösteriminin tanımını da açınca göstermeniz
gereken şudur: $x \in A$ ya da $x \in B$. Böylece ``$A \cup B$'deki her
!!{element}, aynı zamanda $B \cup A$'da da bir !!a{element} olarak bulunur'' sözü
aslında şu anlama gelir: ``her $x$ için, $x \in A$ ya da $x \in B$ ise
$x \in B$ ya da $x \in A$'dır.'' Önermedeki tanımları bütünüyle
açarsak göstermemiz gerekenin şu olduğunu görürüz:

\begin{prop}
Herhangi $A$ ve $B$ kümeleri için: (a) her $x$ için, $x \in A$ ya da
$x \in B$ ise $x \in B$ ya da $x \in A$; ve (b) her $x$ için, $x \in
B$ ya da $x \in A$ ise $x \in A$ ya da $x \in B$.
\end{prop}

Önemli olan, tanımları açmanın bir kanıt kurmanın zorunlu bir parçası
olmasıdır. Bunu doğru yapmak bazen güçtür: tanımdaki değişkenleri ve
kanıtladığınız iddiadaki terimleri dikkatle ayırmalı ve eşleştirmelisiniz.
Başarılı olmak için sorunun ne istediğini ve soruda kullanılan bütün
terimlerin ne anlama geldiğini bilmelisiniz---çoğu kez birden fazla
tanımı açmanız gerekir. Aşağıdakiler gibi basit kanıtlarda çözüm,
tanımların kendilerinden neredeyse hemen çıkar. Elbette her zaman bu
kadar basit olmayacaktır.

\begin{prob}
$A \cap B \neq \emptyset$ olduğunu kanıtlamanız istendiğini varsayın.
Burada geçen bütün tanımları açın; yani ifadeyi ``$\cap$'', ``='' veya
``$\emptyset$'' kullanmadan yeniden söyleyin.
\end{prob}

\end{document}
